\section{Conclusion}

We have derived one-loop matching conditions for non-singlet quark distributions, including unpolarized, helicity and transversity distributions. The matching condition, which can also be viewed as a factorization, connects the quasi quark distribution to the light-cone distribution measurable in experiments, and thereby allows an extraction of the latter from the former, which is defined as a time-independent correlation at finite nucleon momentum and therefore can be evaluated on the lattice. To carry out a lattice simulation of the quasi quark distribution, one employs the transverse momentum cutoff regulator for UV divergences. This choice of UV regulator generates a linear divergence multiplied by a singular factor in the axial gauge $A^z=0$.
The origin of this is the self energy of the Wilson line. In our one-loop factorization formula this divergence can be removed with a principal value prescription. However, to eventually achieve a factorization to all-loop orders, it is important to
investigate how this divergence structure affects the matching between quasi distribution and light-cone distribution beyond one-loop level.

\begin{acknowledgments}
We thank J. W. Chen for useful discussions about the $1/(1-x)^2$ singularity. This work was partially supported by the U.
S. Department of Energy via grants DE-FG02-93ER-40762, and a grant (No. 11DZ2260700) from the Office of
Science and Technology in Shanghai Municipal Government, and by National Science Foundation of China (No. 11175114).
\end{acknowledgments}
